\begin{spacing}{1.2}
	\chapter{PENDAHULUAN}
\end{spacing}

\pagenumbering{arabic}
\vspace{4ex}

\section{Latar Belakang}
Perkembangan teknologi Internet of Things (IoT) memungkinkan berbagai perangkat untuk saling terhubung dan bertukar data melalui jaringan internet. Salah satu implementasi IoT yang banyak digunakan adalah sistem monitoring berbasis kamera yang dapat digunakan untuk mendukung keamanan dan pengawasan suatu lingkungan.

ESP32-CAM merupakan salah satu perangkat yang banyak dimanfaatkan dalam pengembangan sistem monitoring karena memiliki modul kamera dan konektivitas Wi-Fi dalam satu perangkat dengan biaya yang relatif rendah. Namun, sistem monitoring berbasis kamera memerlukan infrastruktur yang mampu mengelola pengiriman data, penyimpanan gambar, pemantauan kondisi perangkat, serta penyajian informasi kepada pengguna secara real-time.

Pada penelitian ini dikembangkan sebuah platform monitoring CCTV berbasis ESP32-CAM yang terintegrasi dengan website dan aplikasi mobile. Sistem memanfaatkan protokol MQTT untuk komunikasi data antar perangkat dan WebSocket untuk menampilkan pembaruan data secara real-time pada website tanpa mekanisme polling. Data hasil tangkapan kamera disimpan pada object storage, sedangkan informasi perangkat dan aktivitas sistem disimpan pada basis data terpusat.

Selain itu, sistem dilengkapi dengan fitur motion detection untuk mendeteksi adanya aktivitas pada area pengawasan, monitoring telemetri perangkat untuk mengetahui kondisi kamera secara real-time, serta notifikasi mobile yang dapat memberikan informasi kepada pengguna ketika terjadi aktivitas yang terdeteksi oleh sistem. Seluruh layanan dijalankan pada server Intel NUC menggunakan teknologi containerisasi sehingga memudahkan proses pengelolaan dan pengembangan sistem.

\section{Rumusan Masalah}
Berdasarkan uraian latar belakang di atas, rumusan masalah dalam penelitian ini didefinisikan secara lebih spesifik sebagai berikut:
\begin{enumerate}
	\item Bagaimana merancang dan mengimplementasikan platform monitoring CCTV berbasis ESP32-CAM yang dapat diakses melalui website dan aplikasi mobile secara real-time?
	\item Bagaimana membangun mekanisme komunikasi data dan penyimpanan terpusat menggunakan MQTT, WebSocket, dan server Intel NUC sehingga sistem dapat berjalan secara stabil dan efisien?
	\item Bagaimana mengimplementasikan fitur motion detection, notifikasi mobile, dan monitoring kondisi perangkat untuk mendukung proses pengawasan dan pengelolaan sistem?
\end{enumerate}

\section{Tujuan Penelitian}
Sejalan dengan rumusan masalah, penelitian ini bertujuan untuk:
\begin{enumerate}
	\item Menghasilkan rancang bangun perangkat keras dan kelistrikan sistem drone otonom yang mampu menyuplai daya komputasi secara stabil bagi perangkat \textit{edge} Raspberry Pi 5 tanpa memicu \textit{voltage sag}.
	\item Mengevaluasi performa iterasi model YOLOv8n format ONNX yang dijalankan pada sistem operasi \textit{headless} untuk meraih keseimbangan \textit{framerate} (FPS) dan latensi terbaik.
	\item Membangun mesin logika otonom tipe \textit{Persistent Multi-Target Scouting} berbasis MAVSDK yang memampukan wahana untuk menavigasi, mendokumentasi korban, dan menghindari pelacakan target duplikat secara mandiri.
\end{enumerate}

\section{Batasan Masalah}
Agar ruang lingkup penelitian tetap terarah dan terukur, ditetapkan batasan masalah sebagai berikut:
\begin{enumerate}
	\item Sistem monitoring menggunakan perangkat ESP32-CAM yang terhubung ke server melalui jaringan internet menggunakan protokol MQTT.
	
	\item Sistem terdiri dari website monitoring real-time berbasis WebSocket dan aplikasi mobile yang digunakan untuk pemantauan serta penerimaan notifikasi.
	
	\item Penyimpanan data menggunakan MySQL untuk metadata dan MinIO Object Storage untuk penyimpanan citra hasil tangkapan kamera.
	
	\item Fitur yang dikembangkan meliputi motion detection, monitoring telemetri perangkat (RSSI, uptime, free heap, dan status perangkat), serta pengiriman notifikasi kepada pengguna.
	
	\item Infrastruktur sistem dijalankan pada server Intel NUC berbasis Debian 12 dengan dukungan Docker, Cloudflared Tunnel, dan Webmin untuk pengelolaan layanan.
\end{enumerate}

\section{Manfaat Penelitian}
Penelitian ini diharapkan dapat memberikan kontribusi dan manfaat sebagai berikut:

\begin{enumerate}
	\item Bagi pengguna dan pengelola sistem, penelitian ini menghasilkan platform monitoring CCTV berbasis IoT yang memungkinkan pemantauan kondisi kamera, penerimaan notifikasi aktivitas secara real-time, serta pengelolaan perangkat dan layanan melalui website maupun aplikasi mobile.
	
	\item Bagi pengembangan teknologi, penelitian ini dapat menjadi referensi dalam pengembangan sistem monitoring berbasis IoT yang memanfaatkan ESP32-CAM, MQTT, WebSocket, motion detection, dan infrastruktur server terpusat untuk mendukung implementasi sistem pengawasan yang lebih efektif dan mudah dikembangkan di masa mendatang.
\end{enumerate}